\chapter{Introduction}
\label{ch:intro}

\section{What this document is}

This is the record of building a \emph{nozzle optimizer}: a program
that does not merely compute the flow through a given rocket
nozzle, but searches for the nozzle shape that produces the most
thrust --- and can certify that what it found satisfies the
mathematical conditions of an optimum.

The target application is the rotating detonation engine (RDE), in
which combustion runs continuously around an annular chamber
instead of sitting still. That makes the design problem harder in a
specific way: the nozzle is handed not one operating condition but
a periodic family of them, thousands of times a second. Everything
this document builds is aimed at that difficulty.

The work is organized in two \emph{bricks}. Brick~1, completed
earlier, is the flow solver: a method-of-characteristics march
written in a differentiable programming language, checked against
closed-form solutions and against an independent Fortran design
code. Brick~2 --- this document --- is the climb from a solver to
an optimizer, step by step, each step carrying its own certified
test suite.

\section{Who this document is written for}

No prior exposure to nozzle design, gasdynamics, the method of
characteristics, optimization theory, or rotating detonation
engines is assumed. The three chapters that follow supply all of
it:

\begin{itemize}[itemsep=2pt]
\item Chapter~\ref{ch:nozzle} --- what a nozzle is and does, what
  the design variables are, why an optimum exists at all, what an
  RDE is and why it creates a cycle, and what a gradient and an
  adjoint buy you.
\item Chapter~\ref{ch:moc} --- how the solver computes: the method
  of characteristics, cells, columns, boundaries, and why sharp
  corners must be handled separately.
\item Chapter~\ref{ch:method} --- how this program decides that
  something is true, and the symbol table
  (Table~\ref{tab:notation}) for the whole document.
\end{itemize}

Readers already fluent in these can start at
Chapter~\ref{ch:score}. Every later chapter opens with a short
plain-language statement of what it adds, and defines any new term
where it first appears.

\section{The two ground rules}

Two rules apply to every test suite in this program, and they are
what makes the results claims rather than demonstrations. They are
stated fully in Chapter~\ref{ch:method}; in brief:

\begin{enumerate}[itemsep=2pt]
\item \textbf{No tolerance is chosen by hand.} Every pass band is
  \emph{derived} --- usually by running the same computation at two
  resolutions and letting the scheme measure its own error. A test
  that passes inside a hand-picked tolerance certifies nothing.
\item \textbf{Every suite must be able to say no.} Alongside the
  positive checks, each suite plants deliberate corruptions that it
  is required to catch. If the corrupted run does not fail, the
  suite is worthless and the step does not count as done.
\end{enumerate}

A consequence worth stating: in this program the bugs the tests
catch are results too. The main chapters present each step as it
now works; everything that went wrong on the way --- every trap and
measured structural lesson, with its diagnosis and fix --- is
collected in Appendix~\ref{ch:appendix}, cross-referenced from the
step it belongs to.

\section{The independent referee}

One further principle shapes the work. Alongside this program there
is GENO~\cite{geno}, an independent Fortran nozzle design code:
different language, different discretization, different
thermodynamics, written by other people. By standing decision it is
\emph{never modified and never differentiated}. It plays two roles:
as a \emph{generator} it produces reference contours and flowfields
that our optimizer did not make, and as a \emph{referee} its fields
are required to satisfy our own equations, and ours its numbers.
Two independently built programs agreeing on one physical problem
is evidence that neither can produce alone.

\section{The road through this document}

Figure~\ref{fig:pipeline} maps the whole of Brick~2: each box is a
chapter, in the order the work was done and certified.

The engine is assembled in this order. Chapter~\ref{ch:score}
builds the score (thrust, computed two independent ways) and the
optimality certificate. Chapter~\ref{ch:firstopt} closes the first
optimization loop and checks it against a case whose answer is
known in closed form. Chapter~\ref{ch:wallinput} inverts the
solver so that the wall shape becomes an \emph{input} --- the
structural change free-form design requires --- and records a
finding about differentiating tabulated gas properties.
Chapter~\ref{ch:referee} brings the independent referee online and
confronts Rao's 1958 hand-derived optimality conditions with the
automatic-differentiation engine, term by term, on a real
thrust-optimized contour.

Chapter~\ref{ch:cycle} adds the cycle layer that makes the engine
an RDE tool, with a theorem-level test the implementation cannot
pass by accident. Chapter~\ref{ch:plug} builds the plug (aerospike)
nozzle march, whose free exhaust boundary is genuinely new physics,
and certifies it against a closed-form solution.
Chapter~\ref{ch:plugcycle} takes the certified plug into the cycle
and measures where a truncated plug should be cut ---
and why that is \emph{not} the design for the average condition.
Chapter~\ref{ch:genotwin} runs the strongest external check the
program owns, our march on the independent code's own designed
spike, and reports one honest unresolved finding.
Chapter~\ref{ch:swirl} extends the march to swirling exhaust and
states plainly what that extension does and does not model.
Chapter~\ref{ch:status} is the ledger: the status of every suite,
what a design run can fix, what remains, and the reproduction
commands.

\begin{figure}[htbp]
\centering
\resizebox{\textwidth}{!}{%
\begin{tikzpicture}[
  box/.style={draw, rounded corners=2pt, minimum height=1.55cm,
    minimum width=2.1cm, align=center, font=\footnotesize,
    text width=1.9cm, thick},
  done/.style={box, draw=stepdone, text=stepdone},
  ahead/.style={box, draw=stepahead, text=stepahead},
  arr/.style={-{Stealth[length=2mm]}, thick, color=black!60},
  node distance=0.32cm]
\node[box, draw=black!55, text=black!55] (b1)
  {diff.\ march\\(Brick 1)};
\node[done, right=of b1] (s1) {score $J$,\\two routes\\8/8};
\node[done, right=of s1] (s2) {gradient +\\corner meter\\5/5};
\node[done, right=of s2] (s3) {optimizer vs\\closed form\\4/4};
\node[done, right=of s3] (s4) {wall as\\INPUT\\5/5 + 2/2};
\node[done, right=of s4] (s5) {referee\\GENO + Rao\\9/9};
\node[done, right=of s5] (s6) {cycle layer\\+ oracles\\5/5, 6/6, 8/8};
\node[done, right=of s6] (s7) {plug march:\\cell, oracle, axi\\6/6, 6/6, 5/5};
\node[done, right=of s7] (s8) {truncated plug\\under cycle\\5/5, 9/9};
\node[done, right=of s8] (s9) {full-field\\plug twin\\8/8};
\node[done, right=of s9] (s10) {free-vortex\\swirl\\9/9};
\node[ahead, right=of s10] (s11) {measured cycle,\\TR-SQP\\(ahead)};
\draw[arr] (b1) -- (s1);
\draw[arr] (s1) -- (s2);
\draw[arr] (s2) -- (s3);
\draw[arr] (s3) -- (s4);
\draw[arr] (s4) -- (s5);
\draw[arr] (s5) -- (s6);
\draw[arr] (s6) -- (s7);
\draw[arr] (s7) -- (s8);
\draw[arr] (s8) -- (s9);
\draw[arr] (s9) -- (s10);
\draw[arr] (s10) -- (s11);
\end{tikzpicture}%
}
\caption{The Brick-2 pipeline. Each box is a chapter of this
document and carries its own certified test suite, with its pass
count shown. Blue: done and certified. Red: the road ahead. Grey:
Brick~1, the prerequisite solver.}
\label{fig:pipeline}
\end{figure}

\section{Environment and provenance}

Work of record ran on host \code{s2} in
\code{RDE\_Design-rde-nozzle-program/} (environment \code{.venv-a1}:
JAX 0.10.2, SciPy, Cantera 3.2.0, Matplotlib). This host's tree is
not a version-controlled checkout: files are ported to the
repository for commit, with a handoff note in \code{validation/}.
The referee code GENO was rebuilt from source on this host
(Chapter~\ref{ch:referee}). Every ``results of record'' block quotes
the carrier script whose output it is, and
Chapter~\ref{ch:status} lists the one-command reproduction for each.
